\begin{cnabstract}
随着低轨（Low Earth Orbit, LEO）卫星通信系统和非地面网络（Non-Terrestrial Networks, NTN）快速发展，手机直连卫星正在从应急通信逐步走向更广泛的直接到终端接入场景，在提升广域接入能力和通信灵活性的同时，也给无线电安全监管带来新的挑战。复杂城市环境中，非合作直连卫星终端可能借助建筑遮挡和复杂传播条件进行隐蔽接入或违规辐射，异常线索通常只能提供粗粒度可疑区域，难以直接确定目标数量和空间位置。相较于传统固定监测站方案，低成本无人机（Unmanned Aerial Vehicle, UAV）平台具有机动性强、部署灵活和空间视角可调等优势，能够为复杂街区中的非合作终端定位排查提供新的技术路径。然而，受城市多径、非视距传播、局部遮挡以及平台载荷与算力约束影响，传统几何定位方法和静态融合策略难以兼顾定位精度、环境适应性和工程可用性。围绕上述问题，本文分别从单目标连续定位和多目标协同鲁棒定位两个层面开展研究。

针对单目标连续定位中多源定位分支可靠性随传播环境和运动状态动态变化、固定权重融合难以稳定适应复杂场景的问题，本文提出了一种基于一致性增强监督权重学习的单目标自适应融合定位方法。该方法首先构建 AoA、PDR 与 RSSI 三路定位结果的位置级仿射融合模型，并利用 RMS 时延扩展、RSSI 方差、平均接收强度以及多源定位结果间的一致性距离构造 6 维权重预测特征。在离线阶段，本文利用真实位置反解每个样本对应的最优融合权重，将其作为监督学习标签；在在线阶段，采用 XGBoost 回归模型根据当前特征直接预测三路融合权重，并完成位置级自适应融合。实验结果表明，在总体测试集上，所提方法的均方根误差和 95\% 分位误差分别为 6.5843 m 和 12.1864 m，相比 AoA-only、PDR-only 和 EKF 等对比方法均取得更低定位误差，并在低信噪比鲁棒性和连续轨迹贴合稳定性方面表现出较好的适应能力。

针对多目标并发场景中目标数量不确定、UAV 空间扫描观测稀疏以及多径诱发伪峰导致传统峰值判定方法易失效的问题，本文进一步提出了一种基于空间信号表征的多目标协同鲁棒定位方法。该方法结合开放街道地图（OpenStreetMap, OSM）、Site Viewer 工具链与射线追踪模型构建复杂城市传播场，并通过 RSSI 栅格映射与无线电地图（Radio Map）重构，将离散空间观测转化为具有区域结构的信号表征。针对多 UAV 协同网络中的随机节点失效问题，本文构建教师—学生知识蒸馏框架，以完整观测条件下的教师模型指导缺失观测条件下的学生模型训练，使学生网络在残缺观测条件下仍能逼近完整网络的判别能力。实验结果表明，所提方法能够有效减缓随机节点失效导致的性能退化，在中高失效比例下保持较稳定的定位性能，提升复杂场景中的协同鲁棒定位能力。

综上，本文建立了面向低成本无人机平台的非合作直连卫星终端定位识别框架。研究结果表明，物理先验、监督权重学习与知识蒸馏等技术的融合，能够提升复杂环境下的定位精度与鲁棒性，可为复杂场景下的无线电安全监管提供有益参考。

\cnkeywords{无人机；卫星通信；融合定位；无线电地图；知识蒸馏}
\end{cnabstract}
